\chapter{Viabilidad y planificación}
\label{chap:planificacion}

\lettrine{X}{xxx},\dots


\section{Análisis de viabilidad}
\label{sec:viabilidad}

{ \color{red} Aspectos a considerar
    \begin{itemize}
      \item Opcional pero recomendable.
      \item Su objetivo es realizar un análisis de la viabilidad del proyecto, validando sus posibles riesgos.
      \item Suele enfocarse descomponiéndolo en diferentes dimensiones: temporal (duración estimada para un TFG), tecnológica (relacionado con tecnologías usadas, normalmente opensource), económica (relacionado con tiempo de alumno en TFG, recursos requeridos que normalmente es HW ya disponible, el SW a realizar y otro SW normalmente opensource), \dots 
    \end{itemize}
}

\section{Planificación}
\label{sec:planificacion}

{ \color{red} Aspectos a considerar
    \begin{itemize}
      \item Enumeración de tareas en las diferentes iteraciones: Diagrama de Gantt
      \item Planificación inicial
    \end{itemize}
}

\section{Seguimiento}
\label{sec:planificacion}

{ \color{red} Aspectos a considerar
    \begin{itemize}
      \item Seguimiento de la planificación, y explicación de posibles retrasos
    \end{itemize}
}

\section{Análisis de costes}
\label{sec:planificacion}

{ \color{red} Aspectos a considerar
    \begin{itemize}
      \item Costes: Directos: personal, HW; Indirectos: luz, agua, ... ; gastos reserva contingencia; importante pro-ratear el precio de HW que no se amortiza en 6 meses, sino que es válido durante 4 o 5 años.
      \item Importante referenciar una fuente de la que se obtiene el coste por horas de determinado perfil
    \end{itemize}
}